% Ad Atticum 14.22 (ad-atticum-14-22) --- English translation
% Translated by Claude (Anthropic) from the Latin (Perseus).
% Style guide: STYLE.md.

\ciceroLetterOpener{CICERO ATTICO salutem}

\ciceroSection{1} Having been told by Pilia that couriers were being sent to you on the Ides, I have at once scratched off this little something. So: first, I wanted you to know that I am leaving here for Arpinum on the sixteenth before the Kalends of June. Send anything you have for me there, then; though I myself shall be with you any moment. For I want, before I come to Rome, to scent out rather more carefully what is to be. Though I am afraid I shall not be at all off the mark in my conjecture. For it is by no means obscure what those men are working at; and indeed my own pupil --- who is dining with me today --- is much in love with the man whom our Brutus wounded. And if you must know (for I have seen it plain), they fear peace. Their underlying premise (\textit{[Greek: hupothesis]}) --- and they make no secret of it --- is this: that a most illustrious man has been murdered, that the whole commonwealth has been thrown into confusion by his fall, that everything he did will be void as soon as we stop being afraid, that his clemency was his ruin, and that, if he had not employed it, nothing of the kind could have happened to him.

\ciceroSection{2} It comes to my mind, however, that, if Pompeius arrives with a strong army --- which is reasonable (\textit{[Greek: eulogon]}) --- there will certainly be war. This prospect and the thought of it upset me. For what was then permitted to you will not now be permitted to us. We rejoiced openly. Then again, they have it always on their lips that we are ingrate. In no way will it now be allowed us, which then was allowed both to you and to many. \textit{One must, then, put on a brave face} (\textit{[Greek: phainoprosop\=eteon]}) \textit{and go} (\textit{[Greek: iteon]}) \textit{into a camp?} A thousand times better to die --- especially at this age of mine. And so the Ides of March do not console me as much as they did before. For they contain a great flaw. Though those young men ``thrust this reproach away amid other glories'' (\textit{[Greek: allois en esthlois tond' ap\=othountai psogon]}). But if you have any better hope --- since you hear more and are present at counsels --- I should be glad if you would write it to me, and at the same time think about what we ought to do regarding my votive embassy. As for me, in these parts I am warned by many not to be in the Senate on the Kalends. For it is said that troops are secretly being mustered for that day, and indeed against the very men who, it seems to me, will be safer anywhere than in the Senate.
